% Related Work

\paragraph{Agent-based models in finance.}
The modern ABM literature in finance traces to the Santa Fe Artificial Stock
Market~\cite{lebaron2000agent} and the Lux--Marchesi model~\cite{lux1999scaling}, which
showed that populations of heterogeneous traders---chartists who follow trends and
fundamentalists who trade toward fair value---can endogenously produce fat-tailed returns
and clustered volatility. Hommes~\cite{hommes2006heterogeneous} surveys extensions
incorporating adaptive expectations and evolutionary strategy selection. A common
limitation is that agents interact globally (mean-field coupling), so the spatial
structure of contagion is absent.

\paragraph{Lattice and spin models.}
Ising-type models map agent states (buy/sell) onto spins and impose nearest-neighbor
coupling on a lattice~\cite{stauffer2000generalization,chowdhury2000statistical}.
These models naturally capture local herding and can exhibit phase transitions between
ordered (bubble/crash) and disordered (efficient) regimes. Kaizoji~\cite{kaizoji2002speculative}
extends the Ising framework to three states (buy/hold/sell) to model speculative
activity. Our approach generalizes the Ising lattice in two ways: (i)~we use
Wolfram totalistic CA rules, which provide a richer and fully enumerable space of
local update functions, and (ii)~we couple the lattice to an external or endogenous
price signal via boundary agents.

\paragraph{Crash mechanisms and critical transitions.}
Sornette~\cite{sornette2017why} argues that large crashes are the culmination of
cooperative herding behavior and can be modeled as critical points. Gabaix et
al.~\cite{gabaix2003theory} derive power-law tails from the size distribution of
large institutional trades. Scheffer et al.~\cite{scheffer2009early} propose generic
early-warning signals---rising autocorrelation and variance---for systems approaching
a tipping point. Our phase diagram (Section~5) exhibits a discontinuous transition
at a critical minority fraction~$f_c$, consistent with a first-order tipping point
rather than the continuous (second-order) transitions typical of Ising models at
criticality.

\paragraph{Cellular automata in complex systems.}
Wolfram~\cite{wolfram2002new} classifies one-dimensional CA into four behavioral
classes (fixed points, periodic, chaotic, complex). Two-dimensional totalistic rules
on a Von~Neumann neighborhood have been studied for pattern formation and
self-organization, but their application to financial markets remains largely
unexplored. Farmer et al.~\cite{farmer2005predictive} demonstrate that even
zero-intelligence agents can reproduce aspects of market microstructure when placed
in a realistic order-book mechanism, motivating our use of simple deterministic rules
as a minimal model of trader behavior.
